\documentclass[10pt,a4paper]{scrartcl}

\usepackage[T1]{fontenc}
\usepackage[utf8]{inputenc}
\usepackage[ngerman]{babel}
\usepackage{lmodern}
\usepackage[a4paper,left=1.9cm,right=1.9cm,top=1.7cm,bottom=1.6cm]{geometry}
\usepackage{xcolor}
\usepackage{enumitem}
\usepackage{tabularx}
\usepackage{microtype}
\usepackage[hidelinks]{hyperref}

\definecolor{tured}{HTML}{C40D1E}
\definecolor{textgray}{HTML}{434343}

\setkomafont{section}{\normalsize\bfseries\color{tured}}
\setlength{\parindent}{0pt}
\setlength{\parskip}{2pt}
\setlist[itemize]{leftmargin=5mm,itemsep=1.5pt,topsep=1pt}
\pagestyle{empty}

\begin{document}

{\Large\bfseries\color{tured} Tätigkeitsübersicht vom 01.07. bis 12.08.2026}\par
\vspace{4pt}
\rule{\textwidth}{1.2pt}
\vspace{5pt}

\begin{tabularx}{\textwidth}{@{}lX@{}}
\textbf{Name:} & Florian Rzepka \\
\textbf{Zeitraum:} & 01.07.2026 bis 12.08.2026 (sechs Wochen) \\
\textbf{Schwerpunkt:} & Abschluss der Dissertation (schriftlicher Teil), wissenschaftliche Publikationen und berufliche Orientierung
\end{tabularx}

\section*{Dissertation und wissenschaftliche Qualifikation}
\begin{itemize}
    \item Fortlaufende Ausarbeitung, Zusammenführung und Überarbeitung der Dissertation im Bereich neuronaler Zustandsbestimmung und eingebetteter Batteriesysteme.
    \item Fertigstellung einer ersten vollständigen Fassung am \textbf{17.07.2026}.
    \item Weiterführende inhaltliche und formale Überarbeitung der vollständigen Fassung bis zum Ende des Berichtszeitraums.
    \item Fortschrittsgespräche mit meiner betreuenden Professorin am \textbf{01.07.2026} und \textbf{21.07.2026}.
    \item Aktuelle Erstellung der Präsentation für die wissenschaftliche Verteidigung. Der Termin befindet sich in Abstimmung und ist voraussichtlich für Ende Oktober 2026 vorgesehen.
\end{itemize}

\section*{Publikationen und wissenschaftliche Zusammenarbeit}
\textbf{Gemeinsames Publikationsvorhaben mit dem Hasso-Plattner-Institut}
\begin{itemize}
    \item Arbeitstreffen mit dem Hasso-Plattner-Institut, Digital Engineering Fakultät, am \textbf{22.07.2026} und \textbf{05.08.2026}.
    \item Inhaltlicher Schwerpunkt: gemeinsames Paper zur Optimierung neuronaler Netze durch Pruning und Quantisierung.
    \item Geplante Einreichung bei der Fachzeitschrift \emph{Engineering Applications of Artificial Intelligence} (EAAI).
\end{itemize}

\textbf{Überarbeitung eines bereits eingereichten EAAI-Manuskripts}
\begin{itemize}
    \item Manuskript: \emph{Embedded Artificial Intelligence in Battery Management Systems: Pruning and Quantization for Efficient State-of-Charge and State-of-Health Estimation}.
    \item Rückmeldung aus dem Begutachtungsverfahren der Fachzeitschrift \href{https://www.sciencedirect.com/journal/engineering-applications-of-artificial-intelligence}{\emph{Engineering Applications of Artificial Intelligence}} am \textbf{12.07.2026}.
    \item Bearbeitung der Gutachterhinweise und erneute Einreichung des überarbeiteten Manuskripts am \textbf{02.08.2026}.
\end{itemize}

\section*{Berufliche Orientierung und Netzwerkaktivitäten}
\begin{itemize}
    \item Unverbindliches Fachgespräch mit einem Mitarbeiter von Solaraero in Köln über Tätigkeiten im Bereich Modellierung von Energiespeichern und die Möglichkeit einer Initiativbewerbung.
    \item Erstes Orientierungsgespräch mit einem ehemaligen Doktoranden der Technischen Universität Berlin über eine mögliche Startup-Gründung zur Optimierung von Energiesystemen. Als mögliche Tätigkeitsfelder wurden Embedded Systems und Modellierung besprochen.
    \item Aktualisierung und fachliche Überarbeitung meines LinkedIn-Profils zur Vorbereitung der weiteren Stellensuche und beruflichen Vernetzung.
\end{itemize}

\vfill
{\small\color{textgray}Berichtszeitraum: 01.07.2026--12.08.2026 \hfill Nächster Bericht: 23.09.2026}

\end{document}
